\section{Vigésima iteración de la Fase 5: Tamaño de la Memoria}
\label{sec:F5I20}

En esta iteración incluimos información sobre el sistema externo, en cuanto a la cantidad de memoria de que disponer, de forma que al generar los cartuchos GWC se indique el porcentaje que van a ocupar, como dato informativo para el usuario. Como cada sistema externo pude disponer de una cantidad de memoria diferente, y el programa final puede ocupar mas o menos de esa memoria, solo damos el tamaño y el porcentaje, y será el usuario quien decide después; la información nunca bloquea.

\subsection{Análisis del estado actual}
\label{sec:F5I20Analisis}

\begin{itemize}
	\item De los dos datos que hacen falta para el porcentaje, el sistema ya
	conoce uno: los bytes del GWC recién creado
	(\texttt{BaseCartucho.generarGWC}, presentados en el panel del Editor por
	\texttt{informarTamanos}). El otro, la memoria de que dispone la máquina
	externa, no se declara actualmente en ningún sitio.
	\item El GWS describe la máquina externa (estilo, resolución, botones,
	teclas, \ldots) pero no dice cuánta memoria tiene el hardware que la
	ejecuta; y desde F5-I19 los cartuchos (GWB/GWI) llevan la información de
	la autoría, con lo cual el sitio natural para declararla son PANTALLAS Y
	CARTUCHOS.
	\item La RAM la ocupa lo que viaja dentro del GWC (sus bytes); la FLASH la
	ocupará el fondo (img\_f), que todavía no se guarda en el cartucho, pero el sistema considerará que si, mirará la ocupación del bloque del fondo, hoy será cero, y dará la ocupación. Cuando guardemos el fondo no hay que tocar nada para que siga funcionando.
\end{itemize}

\subsection{Diseño: el campo \texttt{memoria} (RAM y FLASH)}
\label{sec:F5I20Diseno}

\begin{itemize}
	\item \textbf{Campo opcional \texttt{memoria}, con dos subcampos
		\texttt{RAM} y \texttt{FLASH},} en KB, enteros $\ge 0$, que PUEDEN QUEDAR
	A CERO (o simplemente no venir, lo que vale lo mismo). Se incluye en
	PANTALLAS (GWS) y CARTUCHOS (GWB/GWI), en la raíz del JSON:
	\texttt{"memoria": \{"RAM": 216, "FLASH": 0\}}.
	\item \textbf{¿De dónde sale la memoria vigente al calcular?} De lo que
	MANDA: la memoria del cartucho (GWB/GWI) si la declara; si no, la del
	setup (GWS); si ninguna existe o ambas están a cero, NO se muestra
	ningún porcentaje ni ningún aviso (nada que informar).
	\item \textbf{¿Dónde se informa?} Solo cuando tenemos el tamaño definido, que es al generar el GWC (botón «Generar GWC»), junto a los
	tamaños del par: en la línea de \texttt{informarTamanos} del panel del
	Editor y en el aviso final de \texttt{\#estado} (p.~ej.
	«RAM: 4 de 216 KB (2\,\%) · FLASH: 0 de 216 KB
	(0\,\%)»). El porcentaje es un entero (redondeo clásico);
	si el cálculo da menos de 1\,\%, se muestra «$< 1$\,\%
	» para no leer un 0 engañoso.
	\item \textbf{FLASH:} Tamaño del bloque del fondo, si no existe será 0 KB y 0\,\%, si existe el tamaño del bloque. 
	\item \textbf{Validación:} si un valor declarado no
	es un entero $\ge 0$, aviso y se trata como 0 (es decir, sin
	porcentaje); no bloquea nada, coherente con el carácter informativo de
	toda la iteración.
\end{itemize}

\subsection{Plan de cambios}
\label{sec:F5I20Plan}

\begin{itemize}
	\item \texttt{BaseSetup.js}: en la ficha de Setup, campos
	«RAM (KB)»{} y «FLASH
	(KB)»{} (enteros $\ge 0$, vacío vale 0); los recoge
	\texttt{estado()}; la serialización del GWS solo escribe
	\texttt{memoria} (subcampo a subcampo) con los valores distintos de 0;
	la repoblación muestra la del GWS cargado.
	\item \texttt{BaseFicha.js}: los mismos dos campos en la ficha de
	Cartucho; \texttt{poblar} los lee del GWB/GWI; la serialización los
	escribe en la raíz del JSON (solo los subcampos distintos de 0).
	\item \texttt{BaseCartucho.js}: \texttt{validarSetup} (GWS) y
	\texttt{validarComunes} (cartucho) aceptan \texttt{memoria} si sus dos
	subcampos son enteros $\ge 0$ (o no vienen); valor no válido $\to$
	aviso (F5-I16) y 0. \texttt{\_formatearGWI} y \texttt{\_formatearGWB}
	serializan \texttt{memoria} según lo dicho.
	\item \texttt{BaseDescomponer.js}: en \texttt{generarGWIyGWC}, junto a
	\texttt{informarTamanos}, el cálculo del porcentaje (precedencia
	cartucho $>$ setup) y la línea/aviso de ocupación. Cabeceras de módulos
	al día; A5-anexo, una línea de \texttt{\textbackslash input} nueva.
\end{itemize}

\subsection{Pruebas}
\label{sec:F5I20Pruebas}

\begin{itemize}
	\item Caja nueva C8 (grupo \texttt{gwi\_gwc}, solo consola, al estilo de
	C7), registrada y REFERENCIADA EN F5-I18 (la norma vigente): generar el
	par con un setup cuya \texttt{memoria.RAM} es pequeña y distinta de 0
	$\to$ sale el porcentaje (línea del panel y aviso final); con la
	memoria a cero o sin ella $\to$ sin porcentaje y sin avisos; FLASH
	siempre 0 KB y 0\,\%.
	\item S2 y S5 completas en verde; el cambio no toca el formato GWC ni el
	orden del tick.
\end{itemize}

\subsection{Criterios de aceptación}
\label{sec:F5I20Criterios}

\begin{enumerate}
	\item Un GWS o cartucho sin \texttt{memoria} (o con RAM/FLASH a cero)
	genera SIN porcentajes y SIN avisos de memoria.
	\item Con \texttt{memoria.RAM} declarada, al generar el GWC se muestra el
	porcentaje de RAM junto a los tamaños y en el aviso final; solo
	información, nunca bloqueo.
	\item FLASH queda preparada (campo, cálculo y mensaje) con hoy 0 KB y
	0\,\%: cuando exista el fondo, contará allí sin cambiar el diseño.
	\item RAM y FLASH visibles y editables en las pestañas Setup y Cartucho;
	a cero no se serializan en el JSON.
	\item S2 y S5 en verde al cerrar la iteración.
\end{enumerate}